\chapter{Datos}\label{cap:datos}

\section{Conjuntos de canciones}

La tabla~\ref{tab:conjuntos} resume las canciones que se analizan con la herramienta. Todas se procesan con el mismo
cuaderno y los mismos parámetros, a partir del audio oficial.

\begin{table}[ht]
  \centering
  \caption{Canciones analizadas con la herramienta.}\label{tab:conjuntos}
  \begin{tabularx}{\textwidth}{@{}>{\raggedright}p{6.2cm}c>{\raggedright\arraybackslash}X@{}}
    \toprule
    Conjunto & Tamaño & Uso \\
    \midrule
    Angine de Poitrine, \emph{Vol.~II} & 6 temas & Objeto de estudio (audio oficial) \\
    Referencias corregidas de oído & 4 temas & Evaluación; Yor Zarad, a ciegas \\
    Primus & 3 temas & Control sin cuartos de tono: mide el error \\
    King Gizzard, \emph{Flying Microtonal Banana} (2017) & 9 temas & Rock en 24-EDO (precedente señalado por analistas) \\
    Zappa / King Crimson / Gentle Giant & 6 / 6 / 6 & Rock progresivo comparado \\
    \bottomrule
  \end{tabularx}
\end{table}

En total se midieron 36 canciones: las 6 del álbum, los 3 controles de Primus y 27 de los repertorios comparados. Las
referencias de oído son de canciones del álbum, así que no suman al total.

Las seis canciones del álbum son Mata Zyklek, Fabienk, Sarniezz, Yor Zarad, Utzp y Angor. Los controles de Primus son
«Tommy the Cat», «My Name Is Mud» y «Jerry Was a Race Car Driver». Primus cumple dos papeles: es la comparación propia
del trabajo y, como no usa cuartos de tono, todo cuarto de tono que la herramienta marque en sus canciones es error de
medición. Con él se estima el \emph{piso de ruido} de la herramienta.

\section{Referencias para medir el acierto}\label{sec:referencias}

Para saber si la herramienta acierta se necesitan transcripciones verificadas. Se armaron cuatro:

\begin{itemize}
  \item \textbf{Sarniezz y Fabienk.} Partituras de Guitar Pro y MuseScore, corregidas de oído por el autor. Cada
  corrección se verificó con clips: la pista de guitarra real en un canal y la partitura sintetizada en el otro.
  \item \textbf{Mata Zyklek.} La transcripción de la sesión de KEXP, sin cambios en la primera ronda.
  \item \textbf{Yor Zarad.} La primera referencia armada sin mirar ninguna salida de la herramienta: las alturas salen de
  la tablatura de los PDF de Guitar Pro y los tiempos, de los audios que Guitar Pro exporta por pista. El autor la
  verificó de oído con 20 clips y eligió a ciegas entre versiones en 6 clips más.
\end{itemize}

Además, para Sarniezz, Mata Zyklek y Fabienk se armó una referencia de acordes con \emph{todas} las pistas de la partitura
de Guitar Pro (guitarras I a VI y bajo), ya alineadas al tiempo de la grabación oficial. Solo sirve para evaluar la
armonía (sección~\ref{sec:evaluacion}).

Las partituras publicadas no se usan nunca para guiar el análisis: la herramienta tiene que funcionar con canciones sin
partitura. Utzp era una segunda prueba a ciegas, pero su referencia no pasó la verificación de oído (el ritmo de las
guitarras de acordes) y se usó después como una canción más.

\section{Corpus de estilos}\label{sec:corpus}

La tabla~\ref{tab:corpus} muestra los seis corpus que entran en la comparación final. Cada uno aporta los bloques de
rasgos que tiene.

\begin{table}[ht]
  \centering\small
  \caption{Corpus de estilos usados en la comparación.}\label{tab:corpus}
  \begin{tabularx}{\textwidth}{@{}l>{\raggedright}Xlr@{}}
    \toprule
    Estilo & Corpus & Bloques & Piezas \\
    \midrule
    Rock & Rolling Stone 200 \citep{declercq2011} & armonía, escala, metro & 200 (400 análisis) \\
    Práctica común & Kostka-Payne \citep{temperley2009} & armonía & 46 \\
    Jazz & iRealPro \citep{irealpro} & armonía, metro & 1\,186 \\
    Pop & 9 álbumes de Hooktheory & armonía & 121 \\
    Persa & Radif Corpus \citep{kanani2025} & escala, cuartos de tono & 229 \\
    Indostánica & Hindustani Classical Music Transcription Dataset \citep{singh2022} & escala, cuartos de tono & 50 \\
    \bottomrule
  \end{tabularx}
\end{table}

El corpus de rock contiene las 200 canciones de la lista de la revista Rolling Stone, analizadas por dos analistas (400
análisis). El de práctica común proviene de los ejemplos del libro de armonía de Kostka y Payne y solo tiene armonía. El
de jazz reúne los acordes de los estándares de la aplicación iRealPro. El de pop son 121 temas de 9 álbumes analizados
en Hooktheory (Adele, Reneé Rapp, Luke Combs, Billie Eilish, Taylor Swift, Lil Nas X, Lorde, Ed Sheeran y Tate McRae)
\citep{hooktheory}. El
radif es el repertorio clásico de la música persa, organizado en modos (\emph{dastgah}), con intervalos menores que el
semitono (\emph{koron} y \emph{sori}); la conversión del proyecto da 229 piezas \citep{kanani2025}. El corpus indostánico
trae alturas en 22 \emph{shrutis}, que al llevarse a 24-EDO caen todas en los doce grados regulares.

\paragraph{Descargados pero no usados.} SymbTr (makam turco, unas 2\,200 piezas con alturas en comas de Holder)
\citep{karaosmanoglu2012} y The Session (folk irlandés) se convirtieron a bloques y se retiraron de la comparación
final. El maqam árabe no tiene un corpus simbólico abierto equivalente. ProgGP \citep{loth2023}, el corpus que permitiría
comparar con el rock progresivo, solo se obtiene pidiéndolo a sus autores y no estuvo disponible.

\paragraph{Corpus pequeños.} Solo cuentan para el veredicto los corpus con 100 piezas o más. En un corpus pequeño y
disperso casi todo parece típico; en las figuras, las comparaciones con corpus de menos de 100 piezas (Kostka-Payne e
indostánico) aparecen rayadas.

\section{Fuentes sobre la banda}

Las declaraciones de la banda se documentan con la entrevista de \emph{Noize Magazine} \citep{noize} y el podcast
\emph{Aucard de Tours} \citep{aucard2026}. Del podcast se usó una transcripción de trabajo, no oficial, del tramo
1:33--2:58, en la que Khn describe el dúo de guitarra y bajo microtonales, la construcción por loops, la repetición con
crescendo, la relación estructural con el techno, la tradición del rock instrumental y el uso de 5 y 7 tiempos. Los
análisis externos provienen de los videos de Mathémusique \citep{mathemusique2026}, David Bennett \citep{bennett2026},
Levi McClain \citep{mcclain2026} y Thomas Little \citep{little2019}. La prensa \citep{weller2026,dunlevie2026} se usó
solo para los datos generales de la banda y del instrumento.

\section{Licencias}

El audio de la banda no se distribuye. Los corpus conservan la licencia de su fuente; el repositorio publica agregados,
huellas, scripts y referencias, y no los conjuntos completos cuando la licencia no lo permite. Los pesos de BS-RoFormer
SW no declaran licencia y se usaron solo con fines académicos.
